\subsection{An infinite cuspidal lifting theorem and an actual Selmer obstruction}

\begin{statement}{A.59}{Every rank-one exponent lifts cuspidally centrally}{Proved}
Let $E/\Q$ be an elliptic curve, let $D\subset E$ be a nonempty finite
union of torsion sub-packets, and put $U=E-D$. Write
\[
 \Delta=\pi_1(U_{\overline\Q}),\qquad
 I=\ker\bigl(\Delta\longrightarrow T E\bigr),\qquad
 \Pi^{\rm cc}_U=\pi_1(U)/\overline{[I,\Delta]}.
 \tag{A59.1}
\]
Here a torsion sub-packet means that the differences of its geometric
points are torsion in $E(\overline\Q)$. Every continuous section of
$\pi_1(E)\to G_\Q$ lifts to $\Pi^{\rm cc}_U$. This is the full
profinite quotient, with all coefficient primes retained and with
infinite geometric kernel over $\pi_1(E)$.

In particular, on the least-conductor rank-one cubic of A.43,
\[
 E:\ y^2+y=x^3-1,\qquad D=E[3],\qquad P=(1,0),
 \tag{A59.2}
\]
every $a\in\widehat\Z$ gives a proper Kummer section $s_a$ with class
$a\kappa(P)$, and every $s_a$ lifts to $\Pi^{\rm cc}_U$. The fibre of
section classes over $[s_a]$ is a torsor under
\[
 \widehat{\Q^\times}\ \times\
 \widehat{\Q(\zeta_3)^\times}\ \times\
 \widehat{\Q(\sqrt[3]{3})^\times}^{\,2}.
 \tag{A59.3}
\]
The completion here is $\widehat{L^\times}=\varprojlim_m
L^\times/(L^\times)^m$. The assertion is about abstract global lifts;
it does not impose the Selmer condition. In fact A.60 excludes that
condition for explicit nonintegral exponents.
\end{statement}
\begin{proof}
The published theorem of \citet[Theorem~5.10]{AMBES2015} states that
every section on a smooth projective curve of positive genus over
$\Q$ lifts along the torsor $F_D$ when $D$ is a union of torsion
sub-packets. Their Proposition~3.5 identifies
$\pi_1(F_D)$ with the maximal cuspidally central quotient. Its
definition is exactly (A59.1): the kernel $I$ becomes abelian and
conjugation by the entire geometric group on that kernel becomes
trivial. Thus the theorem applies to every proper section, without
assuming a rational point behind the section or assuming that its
localizations are geometric.

For precision about the coefficient limit, $F_D$ is a torsor under
$T_D=\operatorname{Res}_{D/\Q}\mathbf G_m$. Its obstruction is one
class in the continuous group
\[
 H^2\bigl(\Q,T(T_D)\bigr)
   \simeq T\bigl(H^2(\Q,T_D)\bigr),
 \tag{A59.4}
\]
where the displayed isomorphism follows from their Proposition~4.2
and $H^1(\Q,T_D)=0$ by Shapiro and Hilbert~90. This group is
torsion-free. A multiple of $F_D$ comes from a line-bundle torsor
because the corresponding $T_D/\mathbf G_m$ torsor is torsion.
The pullback of its Chern class along the section vanishes over
$\Q$, by the Brauer-group and neighbourhood argument in their
\S5; hence the original obstruction vanishes in (A59.4).
This is a vanishing statement with the full compact coefficient
module. No passage from unrelated finite-level lifts to a supposedly
compatible lift is being made.

The geometric kernel over $\pi_1(E)$ is
\[
 M_D=T(T_D)=\widehat\Z(1)[D(\overline\Q)],
 \tag{A59.5}
\]
with its permutation action multiplied by the cyclotomic character.
It has rank $\deg D$, not rank $\deg D-1$: the latter is the toric
rank in the geometric abelianization, whereas (A59.5) retains the
central commutator of the two compact elliptic generators.
After choosing one lift of a fixed proper section, all lifts modulo
conjugation by $M_D$ are parametrized by $H^1(\Q,M_D)$. Passing from
a fixed proper representative to its geometric conjugacy class
introduces no further quotient. Indeed its centralizer in $T E$ is
$H^0(\Q,T E)=0$, since the rational torsion group of an elliptic
curve is finite.

In (A59.2) the nine punctures are the nine flexes. Their residue
algebra is
\[
 \Q\ \times\ \Q(\zeta_3)\ \times\
 \Q(\sqrt[3]{3})\ \times\ \Q(\sqrt[3]{3}).
 \tag{A59.6}
\]
For example the three-division polynomial is $3x(x^3-3)$.
The origin is rational; at $x=0$ the two points have
$y=\zeta_3,\zeta_3^2$; and for $x^3=3$ there are the two separate
orbits $y=1$ and $y=-2$. This also verifies directly that the whole
divisor is one torsion packet. Shapiro and the full Kummer
isomorphism applied to (A59.5)--(A59.6) give (A59.3). The needed
$\varprojlim^1$ of $H^0(L,\mu_m)$ vanishes because these are finite
groups; thus the completion in (A59.3) is the continuous $H^1$,
not merely a list of finite coefficient classes.

Every $a\kappa(P)$ is a continuous proper section, since
$\pi_1(E)=T E\rtimes G_\Q$. If $a\notin\Z$, its compact section
cannot be that of a rational point: in the completion of the
finitely generated group $E(\Q)$, the closed cyclic subgroup
$\widehat\Z P$ meets $E(\Q)$ in $\Z P$. Consequently the resulting
cuspidally central sections are abstract nongeometric sections.
This proves no assertion about their local geometricity.
\end{proof}

The precise comparison with task A.12 matters. Since $\Delta'\subset
I$, one has $\gamma_3\Delta\subset\overline{[I,\Delta]}$ and
$\Delta''\subset\overline{[I,\Delta]}$. Therefore (A59.1) is a
quotient of both the intrinsic class-two quotient and the intrinsic
metabelian quotient. For the nine-punctured elliptic curve these
maps are strict: $\Delta$ is free profinite of rank ten, and a
boundary generator $c$ and an independent compact generator $b$
have nonzero commutator $[c,b]$ in the free class-two quotient,
whereas $[c,b]$ is killed in (A59.1). This can be detected in the
finite Heisenberg quotient at any odd prime. Thus A.59 neither
calculates the intrinsic metabelian section set nor lifts sections
through the punctured congruence covers in A.12.

\begin{statement}{A.60}{Finite coefficient erasures are excluded already by abelian Selmerness}{Proved}
Let $E/k$ be an elliptic curve over a number field, let $N\ge2$, put
$D=E[N]$ and $U=E-D$, and let $R\in E(k)$ have infinite order.
For a nonempty finite set $\mathcal P$ of rational primes, let
$e_{\mathcal P}\in\widehat\Z=\prod_p\Z_p$ have component $1$ at
the primes in $\mathcal P$ and component $0$ at all other primes.
The proper section with compact Kummer class
\[
 (1-e_{\mathcal P})\kappa(R)
 \tag{A60.1}
\]
has no Selmer lift to the geometrically abelian quotient of
$\pi_1(U)$. Here Selmer means that at every finite place the section
is the image of a local interior point or of a local cusp section.
Consequently it has no Selmer lift to any larger geometric quotient.

For (A59.2), this excludes every
\[
 a=n(1-e_{\mathcal P}),\qquad
 n\in\Z\setminus\{0\},\quad
 \varnothing\ne\mathcal P\subset\{\hbox{primes}\}\text{ finite}.
 \tag{A60.2}
\]
In particular the exponent with $a_2=0$ and $a_p=1$ for every odd
$p$ fails the Selmer condition already at the geometric abelian
level. Nevertheless it has all the abstract finite-kernel lifts
of A.44 and the full cuspidally central lifts of A.59.
\end{statement}
\begin{proof}
We give the proof over a fixed finite extension $L/k$ splitting
$E[N]$. A hypothesized Selmer lift over $k$ would restrict to one
over $L$. Choose an auxiliary prime $\lambda\nmid N$ and, for every
nonzero $d\in E[N](L)$, choose a function
\[
 f_d\in L(E)^\times,\qquad
 \operatorname{div}(f_d)=N(d)-N(O).
 \tag{A60.3}
\]
Such a function exists because $N[d-O]=0$ in $\operatorname{Pic}^0(E)$.
It is an invertible regular function on $U$. The induced homomorphism
to $\pi_1(\mathbf G_m)$ factors through the geometric abelianization.
Thus a global abelian Selmer lift would give, for every $d$, one
global class in $H^1(L,\widehat\Z(1))$. Reduce this class modulo
$\lambda$ and denote its Kummer representative by
$c_d\in L^\times/(L^\times)^\lambda$.

Set $M=\prod_{p\in\mathcal P}p$. Fix a finite set $S$ of places of
$L$ containing the archimedean places, the places above $N\lambda M$,
the places of bad reduction of $E$, and the supports of representatives
of every $c_d$. Enlarge it to contain the finitely many vertical
components of the divisors of the functions $f_d$ on a smooth model.
Outside $S$ the torsion sections are disjoint and these functions
have exactly the divisor (A60.3) on that model. In particular, if
$Q$ reduces to $d\ne O$, then
$v(f_d(Q))=N\,v(t_d(Q))$, where $t_d$ is a formal parameter at $d$;
if it reduces to $O$, then every $f_d$ has valuation
$-N\,v(t_O(Q))$. If it reduces to neither zero nor pole, its
valuation is zero. Multiplication by a constant in any chosen
normalization merely enlarges this one finite set $S$.

For $v\notin S$, the group $E(L_v)$ is profinite, with an open
pro-$\operatorname{char}(\mathbf F_v)$ subgroup. Its
$\mathcal P$-primary factor is therefore finite. The idempotent
decomposition gives
\[
 R=Q_v+T_v,\qquad
 Q_v=(1-e_{\mathcal P})R,\qquad
 T_v=e_{\mathcal P}R\in E(L_v)[\mathcal P^\infty].
 \tag{A60.4}
\]
In particular $T_v$ is torsion of order prime to the residue
characteristic. The full local Tate Kummer map is injective on
$E(L_v)$: its source is the complete profinite group
$\varprojlim_m E(L_v)/mE(L_v)=E(L_v)$. Hence $Q_v$ is the forced
compact point of the putative local Selmer section. It cannot lie
in $D$, since (A60.4) would then make $R$ torsion in $E(L_v)$ and
therefore in $E(L)$. The local section must consequently be an
interior point section with point $Q_v$; a cusp section would have
compact point in $D$. Its images under all the functions $f_d$
are the ordinary local Kummer classes of $f_d(Q_v)$. Since the
global classes $c_d$ are unramified in valuation outside $S$,
\[
 v(f_d(Q_v))\equiv0\pmod\lambda
 \qquad(v\notin S,\ d\ne O).
 \tag{A60.5}
\]

Consider the infinitely many distinct rational points
$R_j=[M^j]R$, $j\ge0$. We claim that for each fixed $d\ne O$,
\[
 v(f_d(R_j))\equiv0\pmod\lambda
 \qquad\hbox{for every }v\notin S\hbox{ and every }j\ge0.
 \tag{A60.6}
\]
There is something to prove only when $R_j$ reduces into $D$.
Use the primary decomposition of the finite reduction group. The
reduction of $Q_v$ has no $\mathcal P$-primary component, and
multiplication by $M^j$ is invertible on this complementary factor.
Since $[M^j](\overline Q_v+\overline T_v)$ is $N$-torsion, it follows
that $N\overline Q_v=0$ and
$N[M^j]\overline T_v=0$ separately. Reduction on prime-to-residue-
characteristic torsion is injective. Because all the $N$-torsion
is split over $L$, there are unique $d_0,d_1\in E[N](L)$ with
\[
 \overline {d_0}=\overline Q_v,\qquad
 [M^j]T_v=d_1.
\]
The unique torsion point reducing to $\overline {R_j}$ is
$d'= [M^j]d_0+d_1$, and
\[
 R_j-d'=[M^j](Q_v-d_0).
 \tag{A60.7}
\]
Multiplication by the $v$-adic unit $M^j$ is a formal-group
isometry: its linear coefficient is a unit and its remaining
terms have strictly greater valuation on the formal group.
Thus the positive formal distance $r$ of $Q_v$ from $d_0$ is
exactly the distance of $R_j$ from $d'$.
If $d_0\ne O$, apply (A60.5) to $f_{d_0}$; if $d_0=O$, apply it
to any one $f_e$ with $e\ne O$. The valuation is respectively
$Nr$ or $-Nr$, so $\lambda\nmid N$ implies $\lambda\mid r$.
At $R_j$, any fixed $f_d$ now has valuation $0$, $Nr$, or $-Nr$,
which proves (A60.6).

The set
\[
 L(S,\lambda)=\{c\in L^\times/(L^\times)^\lambda:
                   v(c)\equiv0\pmod\lambda\ (v\notin S)\}
 \tag{A60.8}
\]
is finite. Indeed its kernel under the ideal-class map is the
finite group of $S$-units modulo $\lambda$th powers and its image
is contained in the $\lambda$-torsion of the finite $S$-ideal-class
group. Fix one $d\ne O$. By (A60.6), all the values $f_d(R_j)$
have classes in (A60.8), so the infinitely many $R_j$ lift to a
finite collection of twists
\[
 C_c:\quad z^\lambda=c^{-1}f_d,\qquad c\in L(S,\lambda).
 \tag{A60.9}
\]
The smooth projective normalization of each $C_c$ is geometrically
connected: the valuation $N$ at $d$ is not divisible by $\lambda$,
so $f_d$ is not a $\lambda$th power even over $\overline L(E)$.
The cover of $E$ has degree $\lambda$, full ramification at $d$
and $O$, and no other ramification. Riemann--Hurwitz gives
\[
 2g(C_c)-2=2(\lambda-1),\qquad g(C_c)=\lambda\ge2.
\]
By Faltings's theorem, each $C_c(L)$ is finite
\citep[Satz~7, p.~365; with the 1984 correction]{AMFaltings1983,AMFaltings1984Erratum}.
This contradicts the infinitely many distinct $R_j$ lying under
points of the finite family (A60.9), and proves the assertion.
\end{proof}

For the cubic, take $N=3$ and $\lambda=2$, so the final finite
family consists of genus-two double covers. Over
$L=\Q(\zeta_3,\sqrt[3]{3})$, the eight functions (A60.3) can be
given completely explicitly. Put $w^2+w+1=0$ and $b^3=3$. The
three Galois orbits are represented by
\[
 \begin{array}{c|c|c}
 d& f_d&f_d(P)\\\hline
 (0,w)&w-y&w\\
 (b,1)&b^2x-y-2&b^2-2\\
 (b,-2)&1-b^2x-y&1-b^2
 \end{array}
 \tag{A60.10}
\]
and their conjugates. These are the flex tangents divided by the
tangent at $O$, normalized to leading coefficient $+1$ at $O$
for the parameter $-x/y$. Their divisors are exactly $3(d-O)$.
The exact script \texttt{a\_main\_rankone\_units.py} computes the
Hessian and division polynomial, verifies all tangent
intersection multiplicities and normalizations, and checks the
finite primary-component implication used in (A60.7) on 857
cyclic test cases. These finite checks supplement, rather than
replace, the uniform argument above.

The same proof applies to every elliptic twist equipped with its
full flex divisor: after choosing a flex as origin over a finite
extension, that divisor is $E[3]$. Thus it is independent of
conductor, analytic rank, and the CM property. Its input is a
compact Kummer class of the precise form (A60.1) relative to that
origin and a nontorsion rational point over the fixed field. It
does not identify arbitrary affine exponent coordinates after a
change by a nontorsion origin.

The mechanism is already a finite coefficient toric constraint
above the full compact Tate module, with a Selmer condition: one
retains the mod-$\lambda$ characters of all the finitely many
functions (A60.3). The simultaneous image in
$\pi_1(E)\times_{G_L}\prod_{d\ne O}\pi_1(\mathbf G_m)_{\mu_\lambda}$
has finite kernel over $\pi_1(E)$ and is a quotient of the geometric
abelianization. This explains exactly why A.60 is consistent with
A.44: abstract sections of that finite-kernel quotient exist, but
their local values cannot all arise from points or cusps. It also
rules out (A60.2) at every larger quotient in task A.12 which
retains this abelian information. No specific congruence cover is
identified with this finite toric quotient without constructing
that comparison. The proof does not give an effective bound for
an offending place, nor does it settle general
$a\in\widehat\Z\setminus\Z$. For an infinite set $\mathcal P$ the
removed component in (A60.4) need not be finite at all but finitely
many places; the argument as written does not extend to that case.
